% ============================================================
%  preamble.tex —— 《椭圆曲线的同调镜像对称》中译本 导言区
%  编译：XeLaTeX（需要 ctex + 中文字体，Windows 下 fontset=windows）
%  请勿在此文件外使用未定义的宏，也不要自行加载宏包
% ============================================================

% ---------- 中文支持 ----------
\usepackage[UTF8, heading=false, fontset=windows]{ctex}

% ---------- 数学 ----------
\usepackage{amsmath}
\usepackage{amssymb}
\usepackage{amsthm}
\usepackage{mathtools}
\usepackage{mathrsfs}
\usepackage{bm}

% ---------- 排版与图形 ----------
\usepackage[a4paper, margin=2.6cm, top=2.7cm, bottom=2.7cm]{geometry}
\usepackage{graphicx}
\usepackage{xcolor}
\usepackage{enumitem}
\usepackage{booktabs}
\usepackage{array}
\usepackage{longtable}
\usepackage{caption}
\usepackage{float}
\usepackage{listings}

% ---------- 超链接 ----------
\usepackage[colorlinks=true,
            linkcolor=NavyBlue,
            citecolor=NavyBlue,
            urlcolor=NavyBlue,
            bookmarksnumbered=true,
            CJKbookmarks=true]{hyperref}

\definecolor{NavyBlue}{RGB}{0,0,128}
\definecolor{DeepRed}{RGB}{150,0,0}
\definecolor{SlateGray}{RGB}{90,90,90}

% ---------- 定理环境（原文 Definition / Proposition / Lemma / Corollary / Proof） ----------
% 与原文字号风格一致：同章节内所有定理类环境共用一个计数器
\theoremstyle{definition}
\newtheorem{definition}{定义}[section]
\newtheorem{example}[definition]{例}
\newtheorem{remark}[definition]{注}
\newtheorem{construction}[definition]{构造}

\theoremstyle{plain}
\newtheorem{theorem}[definition]{定理}
\newtheorem{proposition}[definition]{命题}
\newtheorem{lemma}[definition]{引理}
\newtheorem{corollary}[definition]{推论}

\renewcommand{\proofname}{\textbf{证明}}

% 无编号版本（原文中不带编号的 “Definition.” “Corollary.” 等整段斜体陈述用它）
\newtheorem*{definition*}{定义}
\newtheorem*{theorem*}{定理}
\newtheorem*{proposition*}{命题}
\newtheorem*{lemma*}{引理}
\newtheorem*{corollary*}{推论}
\newtheorem*{remark*}{注}
\newtheorem*{conjecture*}{猜想}

% ---------- 常用记号 ----------
\newcommand{\ZZ}{\mathbb{Z}}
\newcommand{\CC}{\mathbb{C}}
\newcommand{\RR}{\mathbb{R}}
\newcommand{\QQ}{\mathbb{Q}}
\newcommand{\PP}{\mathbb{P}}
\newcommand{\GG}{\mathbb{G}}
\newcommand{\GGm}{\mathbb{G}_m}
\newcommand{\Ainf}{A_\infty}
\newcommand{\db}{D^{b}}
\newcommand{\Dinf}{D_\infty}
\newcommand{\cO}{\mathcal{O}}
\newcommand{\cF}{\mathcal{F}}
\newcommand{\cL}{\mathcal{L}}
\newcommand{\cA}{\mathcal{A}}
\newcommand{\cB}{\mathcal{B}}
\newcommand{\cC}{\mathcal{C}}
\newcommand{\cD}{\mathcal{D}}
\newcommand{\cE}{\mathcal{E}}
\newcommand{\cQ}{\mathcal{Q}}
\newcommand{\cX}{\mathcal{X}}
\newcommand{\cG}{\mathcal{G}}
\newcommand{\cH}{\mathcal{H}}
\newcommand{\cI}{\mathcal{I}}
\newcommand{\cJ}{\mathcal{J}}
\newcommand{\cK}{\mathcal{K}}
\newcommand{\cM}{\mathcal{M}}
\newcommand{\cN}{\mathcal{N}}
\newcommand{\cP}{\mathcal{P}}
\newcommand{\cR}{\mathcal{R}}
\newcommand{\cS}{\mathcal{S}}
\newcommand{\cT}{\mathcal{T}}
\newcommand{\cU}{\mathcal{U}}
\newcommand{\cV}{\mathcal{V}}
\newcommand{\cW}{\mathcal{W}}
\newcommand{\cY}{\mathcal{Y}}
\newcommand{\cZ}{\mathcal{Z}}

\DeclareMathOperator{\Hom}{Hom}
\DeclareMathOperator{\Ext}{Ext}
\DeclareMathOperator{\End}{End}
\DeclareMathOperator{\Aut}{Aut}
\DeclareMathOperator{\Ob}{Ob}
\DeclareMathOperator{\Mod}{Mod}
\DeclareMathOperator{\Fuk}{Fuk}
\DeclareMathOperator{\Coh}{Coh}
\DeclareMathOperator{\Sym}{Sym}
\DeclareMathOperator{\Spec}{Spec}
\DeclareMathOperator{\Rep}{Rep}
\DeclareMathOperator{\GL}{GL}
\DeclareMathOperator{\SL}{SL}
\DeclareMathOperator{\Tr}{Tr}
\DeclareMathOperator{\coker}{coker}
\DeclareMathOperator{\im}{im}
\DeclareMathOperator{\id}{id}
\DeclareMathOperator{\sgn}{sgn}

% ---------- 本节重点（Key Takeaways） ----------
\newenvironment{keytakeaways}{%
  \par\medskip\noindent
  \rule{\linewidth}{0.6pt}\par\nopagebreak\smallskip\nopagebreak
  \noindent{\heiti 本节重点（Key Takeaways）}\par\nopagebreak\smallskip\nopagebreak
  \begin{itemize}[leftmargin=1.7em, itemsep=2pt, topsep=2pt, parsep=0pt]
}{%
  \end{itemize}%
  \vspace{-1pt}\noindent\rule{\linewidth}{0.6pt}\par\medskip
}

% ---------- 代码清单（原文图 4 的 Python 脚本） ----------
\lstdefinestyle{python}{
  language=Python,
  basicstyle=\ttfamily\footnotesize,
  keywordstyle=\color{NavyBlue},
  commentstyle=\color{SlateGray}\itshape,
  stringstyle=\color{DeepRed},
  columns=fullflexible,
  keepspaces=true,
  breaklines=true,
  showstringspaces=false,
  frame=tb,
  rulecolor=\color{black!35},
  numbers=left,
  numberstyle=\tiny\color{SlateGray},
  numbersep=7pt,
  xleftmargin=2.3em,
  aboveskip=4pt,
  belowskip=4pt,
}
\lstset{style=python}

% ---------- 图表标题 ----------
\captionsetup{font=small, labelsep=quad}
\captionsetup[figure]{name=图}
\captionsetup[table]{name=表}

% ---------- 杂项 ----------
\linespread{1.25}
\setlength{\parindent}{2em}
\setlength{\emergencystretch}{2em}
